\section{Conclusion}
\label{sec:conclusion}

We presented \ours, a data-free method for extracting a linear valence
direction from LLM residual streams using just nine LLM-generated stories
per emotion class. The method uses zero labels, zero gradient updates, and
zero prompt engineering. On seven sentiment benchmarks, two word-valence
lexicons, a continuous VAD regression task, and four languages, the same
nine-story-centroid V-axis produces competitive zero-shot results. Data
efficiency is extreme: one story per class already produces a
better-than-random classifier; fifteen stories per class match a supervised
logistic regression trained on 5{,}000 labeled examples.

The most striking interpretability finding is a sudden ``V-shape'' emergence
pattern: the valence direction is present at layer 1 (embedding-level),
disappears through the middle of the network (a 24-layer dilution zone),
then re-crystallizes abruptly at the final 1--2 transformer blocks.
The final-layer V-axis is \emph{nearly orthogonal} to the middle-layer
V-axis (cosine 0.067), showing that the valence direction is freshly
constructed just before the LM head. This pattern replicates across
Qwen, Pythia, and Bloom model families, suggesting it is a universal
feature of causal language model computation.

Equally striking is the valence/arousal asymmetry: the same method that
recovers valence with Pearson $|r|=0.50$ fails to recover arousal at any
layer in any model tested ($|r|\leq0.16$). This suggests that valence is
encoded as a privileged linear direction that arousal does not share, with
implications for LLM interpretability research.

We also report honest negatives: direct prompting beats our probing on raw
classification accuracy; concept generalization to toxicity fails;
cross-LLM ensembling of V-axes hurts. These define the scope of \ours:
it is valuable for efficiency (\speedup{} faster inference), interpretability
(a single inspectable direction), data efficiency (zero labels), and
continuous scoring, but not for pushing state-of-the-art classification
accuracy.

We hope this work will motivate further study of ``when and where'' LLMs
encode concepts as linear directions, and provide a simple reproducible
baseline for data-free concept probing in future interpretability research.
